\section{確率的プログラミングの応用2;変化点と折線回帰}
\subsection{授業内容}
\subsubsection{科目の中でのこのコマの位置づけ}
\subsubsection{コマ主題細目}
変化点検出は，時系列的なデータの中に異なる二つの平均値を持つ群があることをモデリングする手法である。特にある時点から異なる群に属する，という系列的な意味があることと，変化点があるとすればどのあたりになるかという「変化点の位置的不明確さ」を確率分布で表現し，データから検出するという観点は，確率モデルの表現の自由さとデータとの接合を許す確率プログラミング言語の面白さを味わうには最良の材料である。

まずは混合分布モデルのように，二つの群を分類するモデルを再確認し，その上で時系列的なデータという既有知識から「変化点」という考え方の導入，モデリングへと繋げる。
またデータによっては，一定の点を期に線形モデルの傾きが変わるような表現が可能なものがある。この変化点と回帰分析を融合させた，折線回帰モデルを考えることで，固定的なモデルを超えた柔軟なモデリングが可能であることを理解する。

ただしここで使うデータは時系列的なものであるから，一般的な回帰分析の前提であるサンプルの独立性がない。その意味で不適切なモデルであることに注意し，次回の時系列分析へと繋げる。
\begin{description}
	\item[混合分布モデル] データは可視化することが重要であり，見れば明らかに異なる状態の混合であることがわかる場合がある。具体例とともに可視化を行い，またこれまで学んだ混合分布モデルで表現できることを再確認する。ここで用いるデータは，小杉の体重記録データを用いる。
	\item[変化点検出] データの横軸が時系列的な意味を持つのであれば，時空を超えて二つの群が混合しているというのは不自然な前提である。そこで横軸に時系列的な意味を置くと，ある時点から状態が変化したものとして考えることができる。ここでその時点が「いつ」であるのかは不明であるが，わからないことを確率で表現するのが確率モデルの面白い点である。変化点を確率的パラメータとし，その前後で群が異なるというモデルは，変化点検出のモデリングと言われる。このモデリングはポリグラフ検査など，実践的な場面での利用価値も高い。
		\begin{flushright}
			$\to$\citet{Lee201709}のPp.59--61，\citet{松浦201610}のPp.238-245
		\end{flushright}
	\item[折線回帰] 平均点の位置が変わるだけでなく，変化の傾向が明らかな場合は線形モデルを当てはめることができる。変化点の前後で傾きが変わるような線形モデルは，折線回帰とも呼ばれる。折線回帰モデルの実装については，変化点と回帰モデルを組み合わせたあとで，折れる点を繋げる数学的補正を加えたモデルへと修正する。最後に，説明変数が時点であることから回帰分析の前提として標本の独立性が担保されていない問題を指摘する。
\end{description}
\subsubsection{キーワード}
\begin{itemize}
	\item 混合分布モデル
	\item 変化点
	\item 折線回帰
	\item 時系列分析
\end{itemize}
\subsection{授業情報}
\paragraph{コマの展開方法}
講義/遠隔可/演習
\subsubsection{予習・復習課題}
\paragraph{予習}混合分布モデルの応用になるので，混合分布モデルの基本的な書き方や解析方法について，第\ref{lesson2_26}講を復習しておくことが望ましい。
\paragraph{復習}折線モデルが応用できそうなデータを見つけて，自分なりに実践してみると理解が深まるだろう。特に折れる点が複数あるモデルや，折れる点の数を検出するモデルへと拡張するなど，モデル展開の可能性をかんがえることもできる。